\chapter{Relational Persistence: Identity Without Stasis}

\section{State, Gauge, and Form}

Let $X$ be a raw representation space. Its coordinates may contain distinctions that are physically meaningful, distinctions that are merely representational, and distinctions whose relevance depends on the scientific question. Let $G_{\mathrm{gauge}}$ denote transformations that alter representation without altering the content treated as physical in the current model. The gauge-reduced state space is

\[
Q = X/G_{\mathrm{gauge}}.
\]

Some transformations may be physically real yet irrelevant to the identity criterion under investigation. If $G_{\mathrm{perm}}$ denotes those permitted form-preserving transformations, we may define

\[
\bar Q = Q/G_{\mathrm{perm}}.
\]

The second quotient must not be applied mechanically. Rotation can be irrelevant when classifying intrinsic shape and essential when orientation relative to an external field is the observable. Translation can be ignorable for an isolated pattern and decisive in a spatially heterogeneous environment. The theory therefore requires every quotient to state which differences are removed and why.

Let

\[
\Phi : \bar Q \longrightarrow \mathcal{F}
\]

be a form map into a space $\mathcal{F}$ whose coordinates or invariants express the relationships treated as identity-defining. Two reduced states exhibit exact persistence of form when

\[
\bar q_t \neq \bar q_s
\]

while

\[
\Phi(\bar q_t)=\Phi(\bar q_s).
\]

These two equations separate state identity from form identity. The system has genuinely changed because the points of reduced state space differ. It persists because the relation used to identify its form remains unchanged. The philosophical statement that Being can persist through Becoming is a translation of this distinction, not a replacement for it.

\begin{definition}[Exact persistence]
A trajectory $\bar q(t)$ exhibits exact persistence relative to $\Phi$ on a set of times when distinct reduced states can occur while the form image remains identical.
\end{definition}

Real systems often drift, adapt, grow, and repair themselves, so exact equality is too strict. Let $d_{\bar Q}$ be a metric on the reduced state space and $d_{\mathcal F}$ a metric on form space. We call persistence practical when substantial state change coexists with bounded form change. For chosen thresholds $\delta>0$ and $\varepsilon\ge0$,

\[
d_{\bar Q}(\bar q_t,\bar q_s)\ge \delta
\]

while

\[
d_{\mathcal F}\!\left(\Phi(\bar q_t),\Phi(\bar q_s)\right)\le \varepsilon.
\]

This formulation makes explicit what ordinary language leaves implicit. A biological organism, a long-lived social institution, or a recurrent dynamical pattern need not preserve every microscopic invariant to persist. The appropriate question is which relations must remain within tolerance for the organized process to continue as the same form.

\section{The Nontriviality Requirement}

The theory would become empty if any conserved quantity were permitted to define identity without justification. One could choose a map $\Phi$ so coarse that almost every state in the universe became one form, or so fine that persistence vanished under microscopic fluctuations. The selection of defining relationships must therefore be principled. A form map earns scientific use when its invariants correspond to experimentally stable structure, causal organization, symmetry, dynamically preserved constraints, or a clearly declared identity criterion that matters to prediction and intervention.

This requirement also reveals a limitation of pure equivalence-class thinking. Persistence can depend on causal-historical continuity as well as invariant form. A regenerating organism may preserve no simple microscopic invariant across its lifetime while remaining one causally connected, self-maintaining process. A complete theory of persistence should therefore allow a form map to depend not only on instantaneous state but, where necessary, on trajectory history. One may replace $\Phi(\bar q_t)$ by a functional $\Phi[\bar q_{\le t}]$ when the identity criterion is intrinsically historical. This does not abandon invariance; it relocates the invariant from a snapshot to a structured trajectory.

The choice between state-based and history-based identity is itself an example of descriptive adequacy. If a snapshot language loses the continuity relevant to the phenomenon, a process language is required. Relation remains the common theme because both cases ask which dependencies are preserved when the realization changes.

\section{Potential, Constraint, and Actualization}

Persistence must also be distinguished from mere possibility. Let $\mathrm{Poss}$ denote the set or space of admissible relational possibilities. An actualization rule can be represented by

\[
\mathcal{K}_{\mathrm{act}}(r\mid p,\mathcal{C}),
\]

where $p\in\mathrm{Poss}$ and $\mathcal{C}$ summarizes contextual constraints. In deterministic dynamics the kernel reduces to a delta distribution around the flow $r_\Delta=\varphi_\Delta(p)$. In stochastic dynamics it assigns transition probabilities. Potential therefore becomes scientifically meaningful only when the formalism specifies how constraints restrict which possibilities can occur.

If $\mathcal{A}_j$ is an attractor with basin $\mathcal{B}_j$ and initial conditions are sampled from density $\rho_0$, then

\[
\Pr(\mathcal{A}_j)=\int_{\mathcal{B}_j}\rho_0(p)\,dp.
\]

The equation separates several claims that are often conflated. A state can be mathematically admissible without lying on a realized trajectory. A stationary solution can exist without being stable. A locally stable solution can fail to be globally preferred. A globally preferred state can still be infrequently reached under a given distribution of initial conditions. Accordingly, possibility, trajectory, stationarity, local stability, global preference, and actualization frequency must remain distinct throughout the book.

Constraints provide a second description of the same narrowing process. If

\[
\mathcal{M}=\{x:C_a(x)=0\}
\]

is the constraint manifold and the constraint Jacobian has locally constant rank, then

\[
\dim\mathcal{M}=N-\operatorname{rank}(J_C).
\]

If dynamics selects an attracting subset $\mathcal A$, the hierarchy is

\[
X\supseteq\mathcal M\supseteq\mathcal A.
\]

The ambient state space expresses possibility, the constraint manifold expresses compatibility, and the attractor expresses persistence under the chosen dynamics. None of these is interchangeable with the others.

\section{Coherence and Variation}

Let $\mathscr C$ denote a measure of preservation of the defining relational structure and $\mathscr V$ a measure of residual variation after irrelevant collective motions or gauges have been removed. The characteristic regime of coherent persistence is not $\mathscr V=0$ but

\[
\mathscr C\ \text{high},\qquad \mathscr V>0.
\]

A system may therefore be coherent precisely while it changes. This is the formal basis of the statement that coherence is the persistence of defining relationship through variation. The statement is stronger than a metaphor because the next chapters construct explicit instances in which global motion, internal variation, symmetry content, and form distance can be measured separately.

The broader architecture can now be stated without pretending that every transition is already a theorem. Possibility becomes scientifically relevant when relations distinguish alternatives; distinctions become information relative to an observational channel; information and dynamics interact with constraints; constraints make some patterns stable and others unavailable; symmetry classifies classes of form; symmetry breaking and nonlinear selection choose among possibilities; coherent forms can compose into more complex systems; and adaptive systems can eventually support learning, self-modeling, and perhaps reflective subjectivity. The early stages of this sequence are developed mathematically here. The later stages are a research program. The architecture is useful only if the distinction remains visible.
